\section{固定距离传播的精确相位}
\label{chap:feqo-free-propagation}

一维向前运动的正能电子满足 \(E(p)=\sqrt{m^2c^4+c^2p^2}\)。反解正动量分支，\(p(E)=c^{-1}\sqrt{E^2-m^2c^4}\)。设相互作用后第 \(\ell\) 条边带能量为 \(E_\ell=E_0+\ell\hbar\omega>mc^2\)，振幅为 \(c_\ell\)，且只含向前传播支。下面的相位使用精确色散；将探测面上的各边带写成共同幅度归一化，则还要求 \(\max_\ell|v_\ell-v_0|/v_0\ll1\)，其中 \(v_\ell=\partial_pE[p(E_\ell)]\)。若速度差不可忽略，须保留能量归一化到通量归一化的 \(v_\ell^{-1/2}\) 因子，并以概率流而非单纯 \(|\Psi|^2\) 计算到达时间。在这个窄带条件下，位置 \(z=L\) 的波函数为
\begin{equation}
\Psi(L,t)=\sum_\ell c_\ell
\exp\!\left[\frac{\ii}{\hbar}p(E_\ell)L
-\frac{\ii}{\hbar}E_\ell t\right].
\label{eq:feqo-exact-drift-phase}
\end{equation}
每项模长保持 \(|c_\ell|\)，所以能谱不变；干涉项含 \(p(E_\ell)-p(E_{\ell'})\)，所以时间分布会变。

令 \(v_0=\partial_pE(p_0)\)、\(\tau=t-L/v_0\)，提取共同载波 \(e^{\ii p_0L/\hbar-\ii E_0t/\hbar}\)。留下的包络为
\begin{equation}
\psi(\tau;L)=\sum_\ell c_\ell
e^{-\ii\ell\omega\tau+\ii\Phi_\ell(L)},\qquad
\Phi_\ell(L)=\frac{L}{\hbar}[p(E_\ell)-p(E_0)]-\frac{\ell\omega L}{v_0}.
\label{eq:feqo-drift-envelope}
\end{equation}
这是按距离而非按共同传播时间比较相位；二者若混用，会把群延迟重复计算。

Taylor 定理给 \(p'(E_0)=1/v_0\) 和 \(p''(E_0)=-1/(\gamma_0^3mv_0^3)\)。故 \(\Phi_0=0\)、线性项严格抵消，
\begin{equation}
\Phi_\ell(L)=-\frac{\hbar\omega^2L}{2\gamma_0^3mv_0^3}\ell^2
+R_{3,\ell}^{(p)},\qquad
|R_{3,\ell}^{(p)}|\leq\frac{L|\ell\hbar\omega|^3}{6\hbar}
\max_{E\in\mathcal B}|p'''(E)|.
\label{eq:feqo-sideband-phase-expansion}
\end{equation}
\(\mathcal B\) 是实际占据的能量区间。整体相位不改概率，线性相位只平移时间原点，二次相位才能改变峰形。二阶近似的控制量是 \(\max|R_{3,\ell}^{(p)}|\ll1\)，传播距离越长越需检查它；窄能散本身不是充分理由。

若初始波包还有连续能量宽度 \(\sigma_E\)，每一条边带都带着一个邻近能量分布。上述离散式是这些峰足够窄且能量标签可分离时的近似；严格计算应把 \(\sum_\ell\) 外再加上能量积分，并先保持同一电子内部的相干叠加。

这一积分在高斯例子里可以做完。令入射能量相对 \(E_0\) 的偏移为 \(\epsilon\)，取已归一化的能量振幅
\(f(\epsilon)=(2\pi\sigma_E^2)^{-1/4}e^{-\epsilon^2/(4\sigma_E^2)}\)，从而 \(\int|f(\epsilon)|^2\dd\epsilon=1\)。假设 \(\sigma_E\ll\hbar\omega\)，不同 \(\ell\) 的能量峰可分开标号，则探测面振幅在共同的通量归一化因子之外为
\[
\Psi(L,t)=\sum_\ell c_\ell\int_{-\infty}^{\infty}\dd\epsilon\,
f(\epsilon)
e^{\ii p(E_\ell+\epsilon)L/\hbar-\ii(E_\ell+\epsilon)t/\hbar}.
\]
在前述各 \(v_\ell\) 相近、通量因子可共用的范围内，对每个 \(\ell\) 保留 \(p(E_\ell+\epsilon)=p(E_\ell)+\epsilon/v_\ell+O(\epsilon^2)\)，并要求 \(L\max_\ell|p''(E_\ell)|\sigma_E^2/\hbar\ll1\)。剩下的 Fourier 高斯积分给第 \(\ell\) 条边带的时间振幅正比于
\[
c_\ell e^{\ii p(E_\ell)L/\hbar-\ii E_\ell t/\hbar}
\exp\!\left[-\frac{\sigma_E^2}{\hbar^2}
\left(t-\frac{L}{v_\ell}\right)^2\right].
\]
单条边带的时间强度标准差是 \(\hbar/(2\sigma_E)\)；不同边带的中心又在 \(L/v_\ell\) 到达。若这些中心相差足以分开，即使能量振幅 \(c_\ell\) 原本相干，它们在同一时刻也不能充分重叠，理想的周期聚束图像就失效。若 \(L|p''|\sigma_E^2/\hbar\) 不小，上面的高斯还会因群速度色散展宽，应返回未展开的能量积分。

\begin{exercise}
从 \(p(E)=c^{-1}\sqrt{E^2-m^2c^4}\) 直接求 \(p'(E_0),p''(E_0)\)，核对\cref{eq:feqo-sideband-phase-expansion} 的量纲。
\end{exercise}
